%==============================================================================
% فصل دوم: بررسی متون
% دو بخش اصلی: (۱) مبانی نظری و مفهومی  (۲) مرور مطالعات مرتبط
%==============================================================================
\ChapterStart
\chapter{بررسی متون}
\label{ch:literature}
\section{مقدمه فصل}
\label{sec:ch2-intro}
در این فصل ابتدا مبانی نظری و مفهومی مرتبط با خودمراقبتی، پرفشاری خون، آموزش
بیمار، طراحی آموزشی و خرده‌یادگیری تبیین می‌شود؛ سپس مطالعات داخلی و خارجی
مرتبط مرور و جمع‌بندی می‌گردد تا جایگاه پژوهش حاضر در ادبیات موجود روشن شود.

\section{مبانی نظری و مفهومی}
\label{sec:theoretical}
مبانی نظری پژوهش حاضر در گستره مفاهیمی قرار می‌گیرد که شاکله طرح را شکل می‌دهد؛ از تعریف پرفشاری خون به عنوان یک بیماری مزمن و پیش‌رونده، تا مفهوم خودمراقبتی و مدل‌های رفتاری حاکم بر آن، و سرانجام تا روش‌های آموزشی موردبحث، یعنی آموزش کارگاهی و خرده‌یادگیری، و چارچوب طراحی آموزشی \lr{ADDIE}. در ادامه، هر یک از این مفاهیم به‌تفکیک و بر اساس مستندات علمی موجود تشریح می‌شود.

\subsection{پرفشاری خون}
پرفشاری خون اختلالی مزمن و پیش‌رونده است که با افزایش پایدار فشار سیستولی یا دیاستولی خون در سرخرگ‌ها مشخص می‌شود و به دلیل نبود علائم آشکار در مراحل اولیه، غالباً دیر تشخیص داده می‌شود. این بیماری به «قاتل خاموش» شهرت یافته است، چرا که سال‌ها می‌تواند بدون علامت پیش برود و در عین حال آسیب‌های پیش‌رونده‌ای به اندام‌های هدف وارد سازد. بر اساس گزارش سازمان جهانی بهداشت، پرفشاری خون عامل اصلی بیش از ۷/۵ میلیون مرگ در سال در سراسر جهان است و حدود یک‌سوم بزرگسالان جهان را درگیر می‌کند \cite{who2023hypertension}.
\par
معیار تشخیص پرفشاری خون بر اساس میانگین اندازه‌گیری‌های مکرر تعریف می‌شود. بر اساس راهنمای انجمن قلب آمریکا و کالج قلب آمریکا، فشار خون سیستولی مساوی یا بیش از ۱۳۰ میلی‌متر جیوه یا فشار خون دیاستولی مساوی یا بیش از ۸۰ میلی‌متر جیوه به عنوان پرفشاری خون طبقه‌بندی می‌شود. در راهنماهای اروپایی و نیز راهنمای ایران، آستانه ۱۴۰ بر ۹۰ میلی‌متر جیوه همچنان معیار تشخیصی رایج است. در پژوهش حاضر، بیماران تازه‌تشخیص‌داده‌شده با تأیید متخصص قلب و بر اساس این معیارها وارد مطالعه می‌شوند \cite{whelt2017aha,williams2018esc}.
\par
عوارض پرفشاری خون ناشی از آسیب تدریجی به اندام‌های هدف است. افزایش پایدار فشار خون منجر به هایپرتروفی بطن چپ، کاهش انعطاف‌پذیری عروق، افزایش پس‌بار قلبی و در نهایت بروز نارسایی قلبی می‌شود. در حوزه عروق مغزی، پرفشاری خون عامل اصلی سکته مغزی ایسکمیک و خونریزی‌های مغزی است. کلیه‌ها نیز در اثر آسیب تدریجی آرتریول‌های گلومرولی دچار نارسایی پیش‌رونده می‌شوند و در درازمدت، بیماری کلیوی مزمن رخ می‌دهد. علاوه بر این، پرفشاری خون ریسک حمله قلبی، آمبولی، آنوریسم و اختلالات بینایی را افزایش می‌دهد. این عوارض، پرفشاری خون را به یکی از مهم‌ترین عوامل خطر قابل‌مداخله در بیماری‌های قلبی-عروقی تبدیل کرده است \cite{lim2012gbd}.
\par
عوامل خطر پرفشاری خون به دو دسته قابل‌تغییر و غیرقابل‌تغییر تقسیم می‌شوند. عوامل غیرقابل‌تغییر شامل سن، جنس و پیشینه خانوادگی است؛ شیوع بیماری با افزایش سن بالا می‌رود و در مردان تا سن میانسالی بیشتر از زنان است، هرچند پس از یائسگی این اختلاف کاهش می‌یابد. عوامل قابل‌تغییر شامل رژیم غذایی پرنمک، چاقی، کم‌تحرکی، مصرف الکل، استرس مزمن و اختلال خواب است. رژیم غذایی سرشار از سدیم و چربی‌های اشباع، و مصرف ناکافی پتاسیم، منیزیم و کلسیم، از عوامل تغذیه‌ای شناخته‌شده در بروز پرفشاری خون هستند. این عوامل، چارچوب مداخله‌های آموزشی را تشکیل می‌دهند، چرا که بسیاری از آن‌ها از طریق تغییر رفتار قابل کنترل هستند \cite{whelt2017aha}.
\par
اهمیت پرفشاری خون در نظام سلامت ایران ناشی از شیوع بالای آن و بار اقتصادی-درمانی قابل‌توجه آن است. بر اساس گزارش‌های ملی، شیوع پرفشاری خون در بزرگسالان ایرانی حدود ۳۰ درصد تخمین زده می‌شود و این بیماری یکی از علل اصلی مرگ‌ومیر در کشور است. مراجعه به درمانگاه‌ها و بستری‌های ناشی از عوارض قلبی-عروقی، بخش بزرگی از منابع درمانی را به خود اختصاص می‌دهد. همین بار اقتصادی-درمانی، ضرورت اتخاذ رویکردهای آموزشی مؤثر در راستای توانمندسازی بیماران و ارتقای پایبندی آنان به درمان را آشکار می‌سازد \cite{sarrafzadegan2019iran}.

\subsection{مفهوم خودمراقبتی}
خودمراقبتی یکی از مفاهیم بنیادین در پرستاری و مدیریت بیماری‌های مزمن است و طی دهه‌های گذشته سیر تطوری از یک مفهوم ساده‌انگارانه به یک سازه چندبعدی پیچیده طی کرده است. در تعریف کلاسیک، سازمان جهانی بهداشت خودمراقبتی را به عنوان «فعالیت‌هایی که افراد برای حفظ و ارتقای سلامت خود و پیشگیری از بیماری، و محدود کردن اثرات بیماری انجام می‌دهند» توصیف کرده است. این تعریف، خودمراقبتی را به یک گستره وسیع از رفتارهای سلامتی ارجاع می‌دهد که از پیشگیری اولیه تا مدیریت بیماری را در بر می‌گیرد \cite{who2023hypertension,wilkinson2009selfcare}.
\par
در نگاه تخصصی‌تر، خودمراقبتی فرایندی آگاهانه و هدفمند است که طی آن فرد برای حفظ زندگی، عملکرد سلامت و بهزیستی خود، رفتارهایی را انتخاب و اجرا می‌کند. این فرایند شامل سه مؤلفه اصلی است: مراقبت از خود در شرایط عادی، خودمراقبتی در ارتباط با بیماری، و خودمراقبتی در زمان بدتر شدن علائم. این تفکیک، در مدل رایگل برای بیماری‌های مزمن به یکی از مرجع‌ترین الگوهای خودمراقبتی تبدیل شده است \cite{riegel2012middlerange}.
\par
مدل خودمراقبتی دورا اورم یکی از شناخته‌شده‌ترین چارچوب‌های نظری در این حوزه است. اورم خودمراقبتی را «قالب‌بندی رفتار انسان در ساحت‌های آگاهانه و هدفمند» تعریف می‌کند و سه شرط ضروری برای خودمراقبتی را برمی‌شمارد: نخست، شرط آگاهی و شناخت در مورد نیازها و روش‌های پاسخ؛ دوم، شرط توانمندی و برخورداری از مهارت‌های لازم برای انجام رفتارها؛ و سوم، شرط اراده و انگیزه برای انجام خودمراقبتی. در صورتی که یکی از این شرایط نقض شود، عجز خودمراقبتی رخ می‌دهد و نیاز به حمایت پرستاری ایجاد می‌شود. مدل اورم بستر نظری مناسبی برای پژوهش حاضر فراهم می‌آورد، چرا که آموزش خودمراقبتی، در بدنه خود، تلاشی برای تحقق همین سه شرط از طریق ارتقای دانش، مهارت و انگیزه بیماران است \cite{orem1995nursing}.
\par
مدل میان‌ساختاری رایگل، که برای بیماری‌های مزمن توسعه یافته، خودمراقبتی را فرایندی شامل سه بُعد توصیف می‌کند: نگهداری از خود، مدیریت خود و اطمینان از خود. نگهداری از خود شامل رفتارهای پیشگیرانه برای حفظ ثبات وضعیت بالینی است، مانند مصرف منظم دارو، رعایت رژیم غذایی و فعالیت بدنی. مدیریت خود شامل واکنش آگاهانه به تغییرات وضعیت، مانند تنظیم دارو یا مراجعه به پزشک هنگام بروز علائم هشدار است. اطمینان از خود نیز به ارزیابی فرد از اثربخشی رفتارهای خودمراقبتی خود و اعتماد به توان مدیریت بیماری اشاره دارد. این سه بُعد در پرسشنامه خودمراقبتی پرفشاری خون نیز به نوعی بازتاب یافته‌اند \cite{riegel2012middlerange}.
\par
در بیماران مبتلا به پرفشاری خون، رفتارهای خودمراقبتی چندبعدی هستند و شامل پایش منظم فشار خون در منزل با دستگاه دیجیتال، مصرف دقیق داروها بر اساس دستورالعمل پزشک، اتخاذ رژیم غذایی کم‌نمک و کم‌چرب مانند رژیم \lr{DASH}، انجام فعالیت بدنی منظم با شدت متوسط، مدیریت استرس از راهبردهای آرام‌سازی و مدیریت وزن می‌شوند. پایش فشار خون خانگی یکی از مؤلفه‌های کلیدی خودمراقبتی پرفشاری خون است، چرا که تنوع روزانه فشار خون را نمایان می‌سازد و تعامل درمانی بین بیمار و ارائه‌دهنده خدمات سلامت را تسهیل می‌کند. آموزش صحیح نحوه اندازه‌گیری، موقع اندازه‌گیری و تفسیر نتایج، شرط لازم برای اعتبار این پایش است \cite{appel2006dash,han2014hscp}.
\par
عوامل مؤثر بر رفتارهای خودمراقبتی در بیماران مبتلا به پرفشاری خون شامل متغیرهای فردی، اجتماعی و سیستمی هستند. در سطح فردی، دانش در مورد بیماری و درمان، خودکارآمدی، باورهای سلامتی و انگیزه از پیش‌بینی‌کننده‌های مهم رفتار خودمراقبتی هستند. در سطح اجتماعی، حمایت خانواده، ارتباط با تیم درمانی و دسترسی به منابع آموزشی قابل‌اعتماد نقش کلیدی دارند. در سطح سیستمی، طراحی خدمات درمانی به گونه‌ای که آموزش را به روتین پیگیری بیماران متصل کند، ضروری است. پژوهش حاضر بر روی متغیر آموزشی به عنوان یک عامل قابل‌مداخله متمرکز است و بر این پیش‌فرض استوار است که روش آموزش می‌تواند بر دانش، خودکارآمدی و در نهایت رفتار خودمراقبتی بیماران اثر بگذارد.

\subsection{مدل‌های رفتاری مرتبط با آموزش خودمراقبتی}
مدل‌های رفتاری متعددی برای تبیین و پیش‌بینی رفتارهای سلامتی توسعه یافته‌اند و در طراحی مداخله‌های آموزشی به کار گرفته می‌شوند. یکی از این مدل‌ها، مدل باورهای سلامتی است. این مدل بر این فرض استوار است که رفتار سلامتی ناشی از ادراک فرد از حساسیت به بیماری، شدت بیماری، منافع رفتار، موانع رفتار، محرکان عمل و خودکارآمدی است. در حوزه پرفشاری خون، بیمار تنها در صورتی به رفتارهای خودمراقبتی روی می‌آورد که خود را در معرض ریسک بیماری بداند، پیامدهای آن را جدی بگیرد، منافع رفتار را بر موانع آن غالب و ترجیح بیابد و محرکی بیرونی یا درونی برای عمل داشته باشد. آموزش مبتنی بر مدل باورهای سلامتی می‌تواند با بازسازی این باورها به تغییر رفتار کمک کند \cite{rosenstock1974hbm,ramazankhani2019hbm}.
\par
دومین چارچوب، مدل \lr{PRECEDE-PROCEED} است. این مدل یک چارچوب برنامه‌ریزی آموزشی است که بر اساس سه فاز تشخیص، آموزش و ارزیابی تنظیم شده است. در فاز تشخیص، عوامل اجتماعی، اپیدمیولوژیک، رفتاری و محیطی شناسایی می‌شوند. در فاز آموزش، عوامل پیش‌بینی‌کننده رفتار شامل عوامل مساعدکننده، توان‌بخش و تقویت‌کننده مداخله می‌شوند. در فاز ارزیابی، فرایند و پیامد مداخله سنجیده می‌شود. مطالعات داخلی متعدد در حوزه پرفشاری خون بر اساس این مدل طراحی شده‌اند و نتایج آن‌ها نشان‌دهنده اثربخشی آموزش مبتنی بر این چارچوب بر رفتارهای خودمراقبتی بوده است \cite{greenkreuter2005precede,rakhshani2024precede,abedi2020precede}.
\par
سومین مدل، نظریه رفتار برنامه‌ریزی‌شده است. این نظریه رفتار را تابع سه متغیر نگرش به رفتار، هنجارهای ذهنی و کنترل رفتاری ادراک‌شده می‌داند. نگرش به رفتار شامل ارزیابی مثبت یا منفی فرد از انجام رفتار است؛ هنجارهای ذهنی شامل ادراک فرد از انتظارات افراد مهم پیرامون انجام رفتار است؛ و کنترل رفتاری ادراک‌شده شامل باور فرد به توان انجام رفتار است. خودکارآمدی یکی از اجزای کلیدی کنترل رفتاری ادراک‌شده است و آموزش‌های خودمراقبتی غالباً با هدف ارتقای خودکارآمدی بیماران طراحی می‌شوند \cite{ajzen1991tpb}.
\par
چهارمین چارچوب، مدل سبک‌یادگیری \lr{VARK} است. این مدل یادگیرندگان را بر اساس ترجیح حسی-بصری، شنیداری، خواندن-نوشتن و حرکتی-لمسی طبقه‌بندی می‌کند و بر این فرض استوار است که تطبیق روش آموزش با سبک یادگیری فرد می‌تواند اثربخشی آموزش را افزایش دهد. کاربست این مدل در پرفشاری خون نشان داده است که آموزش مبتنی بر سبک یادگیری اختصاصی می‌تواند نمرات سبک زندگی بیماران را به طور معناداری بهبود بخشد. نکته چندبُعدی مهم این است که چارچوب موفق، نه لزوماً پیچیده، می‌تواند در طراحی مداخله تأثیرگذار باشد \cite{fleming2001vark,ghorbani2021vark}.
\par
این مدل‌ها، با تفاوت‌های جزئی، بر یک اصل مشترک تأکید دارند: تغییر رفتار سلامتی، فرایندی چندبُعدی است که از راه دانش، نگرش، خودکارآمدی و محرک‌های محیطی شکل می‌گیرد. آموزش، یکی از راه‌های مداخله بر این مؤلفه‌ها است و طراحی آن باید چارچوب‌مند باشد تا پوشش همه‌جانبه بر متغیرهای مؤثر داشته باشد. در پژوهش حاضر، محتوای آموزشی بر اساس این مبانی تدوین شده و اثربخشی دو روش آموزش بر عملکرد خودمراقبتی بیماران مقایسه می‌شود.

\subsection{آموزش کارگاهی}
آموزش کارگاهی روشی است که در آن گروه کوچکی از یادگیرندگان در یک فضای تعاملی با هدایت مربی، طی یک دوره مشخص، حول موضوع مشخص آموزش می‌بینند. این روش از دهه‌ها پیش به عنوان یکی از مؤثرترین روش‌های آموزش بزرگسالان شناخته شده است و در آموزش‌های سلامت، به‌ویژه آموزش بیماران مزمن، جایگاه ویژه‌ای دارد. نظریه یادگیری بزرگسالان بر تأکید بر تجربه پیشین یادگیرنده، ارتباط محتوا با نیازهای عملی، فعال بودن یادگیرنده و بازخورد فوری استوار است و همه این اصول در آموزش کارگاهی قابل تحقق هستند \cite{knowles1980andragogy}.
\par
ویژگی‌های آموزش کارگاهی شامل تعامل چهره‌به‌چهره، بحث گروهی، تمرین عملی، پرسش و پاسخ و امکان بازخورد فوری است. این ویژگی‌ها، آموزش کارگاهی را به روشی تعاملی و تجربی تبدیل می‌کند که در آن یادگیرنده نقش فعالی دارد و برخلاف سخنرانی یک‌سویه، فقط دریافت‌کننده اطلاعات نیست. در کارگاه، یادگیرنده فرصت می‌یابد تا دانش پیشین، نگرانش و تجربه شخصی خود را با محتوای آموزشی ادغام کند و از هم‌سالان خود بیاموزد. این ویژگی تعاملی در مطالعات متعدد به عنوان تبدیل‌کننده آموزش اطلاعات‌محور به آموزش رفتارمحور توصیف شده است.
\par
در آموزش بیماران مبتلا به پرفشاری خون، آموزش کارگاهی به چند دلیل اهمیت دارد. نخست، بیماری مزمن است و نیازمند یادگیری رفتارهای خودمراقبتی است که باید در زندگی روزمره بیماران ادغام شوند. دوم، بیماران غالباً باورهای نادرست یا نگرانی‌های متعددی در مورد بیماری و درمان دارند که در فضای کارگاهی قابل طرح و بررسی است. سوم، عملکردهایی مانند اندازه‌گیری فشار خون، تفسیر نتایج و تنظیم رفتار، مهارت‌های عملی هستند که با تمرین در کارگاه قابل کسب هستند. چهارم، حمایت هم‌گروهی و تجربه مشترک می‌تواند تقویت‌کننده انگیزه و خودکارآمدی بیماران باشد.
\par
مطالعات متعدد داخلی و خارجی در حوزه پرفشاری خون از آموزش کارگاهی استفاده کرده‌اند. در مطالعات داخلی، آموزش مبتنی بر مدل‌های رفتاری در جلسات کارگاهی منجر به بهبود معنادار نمرات دانش، نگرش و رفتارهای خودمراقبتی بیماران شده است. در مطالعات خارجی نیز مداخله‌های آموزش گروهی مبتنی بر مراقبت حمایتی به بهبود کیفیت زندگی و کنترل فشار خون منجر شده است. این مطالعات به اجماع به این نتیجه رسیده‌اند که آموزش کارگاهی، به‌ویژه زمانی که تعاملی و با پیگیری همراه باشد، می‌تواند به بهبود نتایج بالینی پرفشاری خون کمک کند \cite{rakhshani2024precede,abedi2020precede,nasirian2019group,suprayitno2023telenursing,desouza2016flashcard,kurt2022selfmgmt}.
\par
با وجود این، آموزش کارگاهی با محدودیت‌هایی روبه‌رو است. نخستین محدودیت، زمان و مکان است؛ بیماران برای شرکت در کارگاه باید در زمان و مکان مشخصی حضور یابند که این ممکن است با ساعات کاری، مسافت و مسئولیت‌های خانوادگی تداخل داشته باشد. دومین محدودیت، تداومکی است؛ آموزش در جلسات محدود، در طول زمان تحلیل می‌رود و نیازمند مرور و تقویت است. سومین محدودیت، یکنواختی مخاطبان است؛ کارگاه غالباً با یک محتوا برای همه بیماران طراحی می‌شود، در حالی که نیازها و توانایی‌های فردی متفاوت است. این محدودیت‌ها، زمینه‌ای فراهم می‌آورند که نوآوری‌های آموزشی مانند خرده‌یادگیری به عنوان مکمل کارگاه پیشنهاد می‌شوند.

\subsection{خرده‌یادگیری (\lr{Microlearning})}
خرده‌یادگیری رویکرد آموزشی نوینی است که محتوای آموزشی را در قالب واحدهای کوتاه، متمرکز و قابل‌هضم ارائه می‌کند. طول هر واحد معمولاً چند دقیقه است و محتوا به‌گونه‌ای طراحی می‌شود که بر یک مفهوم یا مهارت مشخص متمرکز باشد. این رویکرد، که در دهه‌های اخیر با گسترش فناوری‌های موبایل و پیام‌رسان‌ها رونق یافته، پاسخی به محدودیت‌های زمانی، حواس‌پرتی‌های محیطی و کاهش دقت و توجه در یادگیرندگان است \cite{monib2024microlearning}.
\par
اصول طراحی خرده‌یادگیری شامل چند اصل کلیدی است. نخست، تمرکز بر یک موضوع مشخص؛ هر محتوای خرده‌یادگیری صرفاً یک نقطه آموزشی را پوشش می‌دهد. دوم، کوتاه بودن؛ محتوا به گونه‌ای طراحی می‌شود که در کمتر از پنج دقیقه قابل تکمیل باشد. سوم، تنوع قالب‌های ارائه؛ خرده‌یادگیری می‌تواند به شکل ویدیو، تصویر، انیمیشن، پادکست، کلیپ متنی یا آزمون کوتاه ارائه شود. چهارم، دسترسی‌پذیری؛ محتوا باید قابل مشاهده در دستگاه‌های موبایل و در هر زمان و مکان باشد. پنجم، ماهیت تکرارپذیر؛ محتوا باید قابل مشاهده مکرر و قابل مرور باشد.
\par
در حوزه سلامت، خرده‌یادگیری به چند دلیل دارای پتانسیل است. نخست، بیماران مزمن نیازمند یادگیری مداوم و مرور هستند؛ خرده‌یادگیری امکان تقویت مستمر و تدریجی را فراهم می‌سازد. دوم، بیماران غالباً محدودیت زمانی دارند و توان شرکت در جلسات آموزشی طولانی را ندارند؛ خرده‌یادگیری با ارائه محتوا در بُعد زمانی کوچک، این محدودیت را کاهش می‌دهد. سوم، محتوای آموزشی می‌تواند با نیازهای شخصی تطبیق یابد؛ مطالعات نشان داده‌اند که محتوای خرده‌یادگیری می‌تواند با وضعیت بالینی و اجتماعی-اقتصادی بیماران شخصی‌سازی شود. چهارم، خرده‌یادگیری می‌تواند با فرایند روزانه بیمار ادغام شود؛ بیمار می‌تواند در زمان‌های پراکنده مانند انتظار در مطب، سفرهای کوتاه یا زمان استراحت، محتوا را ببیند \cite{poorcheraghi2025microlearning,susman2024micropractice}.
\par
شواهد علمی پیرامون اثربخشی خرده‌یادگیری نشان می‌دهد که این روش به‌طور قابل‌اعتمادتر بر شناخت نسبت به رفتار اثر می‌گذارد. یک مرور نظام‌مند در سال ۲۰۲۰ نشان داد که در ۹۱ درصد مطالعات، خرده‌یادگیری منجر به بهبود معنادار پیشاوردهای شناختی شد، در حالی که تنها در ۳۶ درصد مطالعات، تغییر معنادار رفتار رخ داد. این گسست مهم نشان می‌دهد که خرده‌یادگیری به تنهایی برای تغییر رفتاری کافی نیست و باید با مکانیسم‌های پایش، حمایت و تمرین مکرر همراه شود \cite{wang2020microlearning}.
\par
در پژوهش حاضر، خرده‌یادگیری به عنوان مکمل آموزش کارگاهی به کار گرفته می‌شود. هشت ویدیوی آموزشی سه‌دقیقه‌ای در طول چهار هفته به بیماران ارسال می‌شود. این ویدیوها بر محور رفتارهای کلیدی خودمراقبتی پرفشاری خون طراحی شده‌اند: آشنایی با بیماری و روش پایش، تغذیه، فعالیت بدنی، استرس و پایبندی دارویی. ساختار خرده‌یادگیری در این مطالعه با اصول طراحی خرده‌یادگیری همخوان است؛ هر ویدیو بر یک موضوع متمرکز است، سه دقیقه طول می‌کشد و در قالب ویدیویی قابل مشاهده در موبایل بیماران است.

\subsection{الگوی طراحی آموزشی \lr{ADDIE}}
مدل \lr{ADDIE} چارچوب طراحی آموزشی است که در دهه ۱۹۷۰ توسعه یافته و امروزه به یکی از پرکاربردترین مدل‌های طراحی آموزشی در حوزه‌های متنوع، از جمله آموزش سلامت، تبدیل شده است. نام \lr{ADDIE} برگرفته از حروف اول پنج مرحله این مدل است: تجزیه و تحلیل، طراحی، توسعه، اجرا و ارزیابی. این مدل یک فرایند خطی و چرخه‌ای است که هر مرحله به مرحله پیشین وابسته و از خروجی‌های مرحله پیشین استفاده می‌کند \cite{branch2009addie}.
\par
مرحله تجزیه و تحلیل شامل شناخت مخاطبان، موقعیت آموزشی و نیازهای آموزشی است. در پژوهش حاضر، این مرحله شامل بررسی ویژگی‌های دموگرافیک و بالینی بیماران، تحلیل وضعیت خودمراقبتی و شناسایی شکاف‌های دانشی، نگرشی و مهارتی است. خروجی این مرحله، تعریف اهداف آموزشی و تحلیل یادگیرنده، موقعیت و وظیفه است.
\par
مرحله طراحی شامل تدوین اهداف رفتاری قابل‌سنجش، طراحی ساختار دوره و انتخاب روش‌های آموزشی است. خروجی این مرحله، یک نقشه دوره آموزشی است که شامل سرفصل جلسات، محتوای هر جلسه، روش‌های ارائه و ابزارهای ارزیابی است. در پژوهش حاضر، این مرحله شامل طراحی ساختار چهار جلسه کارگاهی و هشت ویدیوی خرده‌یادگیری، تعیین روش‌های سخنرانی، پرسش و پاسخ، بحث گروهی و تمرین عملی است.
\par
مرحله توسعه شامل تولید محتوای واقعی دوره است؛ شامل تهیه اسلایدها، تدوین جزوه و بروشور، تولید ویدیوها و طراحی ابزارهای سنجش. در پژوهش حاضر، این مرحله شامل تولید محتوای آموزشی برای جلسات کارگاهی و ویدیوهای خرد است. ویدیوها هر کدام سه دقیقه و با کیفیت مناسب برای مشاهده در تلفن همراه تولید می‌شوند و بر اساس سرفصل‌های از پیش تعیین‌شده، یک نکته کلیدی را شرح می‌دهند.
\par
مرحله اجرا شامل برگزاری جلسات و ارسال محتوای دیجیتال است. در پژوهش حاضر، جلسات چهارگانه هفتگی در کلاس‌های آموزشی بیمارستان برگزار می‌شود و ویدیوهای خرده‌یادگیری دو بار در هفته از طریق پیام‌رسان به بیماران ارسال می‌شود. پایش حضور و غیاب، مشاهده ویدیوها و بازخورد بیماران در این مرحله ثبت می‌شود.
\par
مرحله ارزیابی شامل ارزیابی فرایندی و ارزیابی پیامد است. ارزیابی فرایندی در طول اجرای مداخله انجام می‌شود و هدف آن اصلاح مداخله است. ارزیابی پیامد در زمان‌های مشخص پس از پایان مداخله انجام می‌شود و هدف آن سنجش اثر مداخله بر پیشاوردها است. در پژوهش حاضر، ارزیابی پیامد شامل سنجش نمره پرسشنامه خودمراقبتی پرفشاری خون و فشار خون در دو نوبت، پیش از مداخله و یک ماه پس از پایان مداخله، است.
\par
مدل \lr{ADDIE} در این پژوهش به عنوان چارچوب ساختاری طراحی دوره انتخاب شده است، چرا که این مدل فرایندی روش‌مند را برای طراحی آموزشی فراهم می‌آورد، که هر مرحله خروجی مشخصی دارد و ارزیابی در تمام مراحل نهادینه شده است. این چارچوب، احتمال تولید دوره‌ای منطبق با نیازهای بیماران را افزایش می‌دهد و روش‌شناختی تحلیل را شفاف می‌سازد.

\subsection{آموزش ترکیبی}
آموزش ترکیبی یا تلفیقی رویکردی است که در آن روش‌های رو‌به‌رو و دیجیتال به گونه‌ای ادغام می‌شوند که نقاط قوت هر یک با نقاط ضعف دیگری جبران شود. در آموزش ترکیبی، جلسات حضوری نقش ایجاد ارتباط، بحث و تمرین عملی را بر عهده دارند، در حالی که محتوای دیجیتال امکان تکرار، مرور و دسترسی در زمان و مکان دلخواه را فراهم می‌سازد. این رویکرد، به جای جانشینی بین روش‌ها، از کنش هم‌افزایی آن‌ها بهره می‌برد \cite{garrison2008blended}.
\par
مطالعات متعدد اثربخشی آموزش ترکیبی را در حوزه بیماری‌های مزمن ارزیابی کرده‌اند. در یک کارآزمایی تصادفی‌شده در بیماران مبتلا به دیابت نوع ۲، آموزش ترکیبی شامل کارگاهسه‌ساعته و دوره آموزشی مبتنی بر وب‌سایت، منجر به بهبود رفتارهای خودمراقبتی در همه ابعاد و افزایش دانش دیابت در مقایسه با آموزش روتین با بروشور شد. در مطالعه‌ای دیگر، خرده‌یادگیری مبتنی بر رسانه‌های اجتماعی، دانش، خودکارآمدی و رفتارهای خودمراقبتی بیماران دیابتی را در مقایسه با کارگاه سه‌ساعته بهبود بخشید. این مطالعات نشان می‌دهند که آموزش ترکیبی می‌تواند بر پیشاوردهای شناختی، رفتاری و خودکارآمدی اثرگذار باشد \cite{ghiyasvandian2023blended,rahbar2024socialmedia}.
\par
با این حال، پژوهش‌های موجود بر چند محدودیت تأکید دارند. نخست، اثربخشی آموزش ترکیبی بر پیشاوردهای بالینی سخت‌تر مانند کنترل گلیسمی یا پیروی دارویی، ناسازگارتر از اثر بر پیشاوردهای نزدیک‌تر مانند دانش و خودکارآمدی است. دوم، اثر افزودن خرده‌یادگیری به کارگاه، در مقایسه با کارگاه به تنهایی، کم‌بودجه است و مطالعات رو‌در‌رو مستقیم نادرند. سوم، کیفیت گزارش‌دهی در مطالعات ترکیبی غالباً ضعیف است؛ جزئیات دوز خرد، تنظیم دقیق پیام‌ها، الگوریتم‌های شخصی‌سازی و دوز به‌طور کامل گزارش نمی‌شوند. این محدودیت‌ها، فضایی برای پژوهش حاضر فراهم می‌آورد که اثر افزودن خرده‌یادگیری به کارگاه را به‌طور مستقیم و در جمعیت بیماران مبتلا به پرفشاری خون ارزیابی کند \cite{lee2022chronicmeta}.


\subsection{چارچوب مفهومی پژوهش}
پژوهش حاضر بر یک چارچوب مفهومی ترکیبی استوار است که چند نظریه و شواهد را ادغام می‌کند. نخست، پرفشاری خون به عنوان یک بیماری مزمن نیازمند خودمراقبتی مداوم است. دوم، خودمراقبتی رفتاری چندبعدی است که بر اساس مدل رایگل شامل نگهداری از خود، مدیریت خود و اطمینان از خود است. سوم، آموزش می‌تواند از طریق ارتقای دانش، خودکارآمدی و انگیزه بر این رفتارها اثر گذارد. چهارم، آموزش کارگاهی روش آموزشی مؤثری است که تعامل، تمرین و بحث را میسر می‌سازد. پنجم، خرده‌یادگیری توان مکمل کردن کارگاه را از راه تقویت، تکرار و دسترسی در زمان دلخواه دارد. ششم، مدل \lr{ADDIE} چارچوبی برای طراحی دوره آموزشی فراهم می‌آورد. هفتم، اثربخشی مداخله با استفاده از پرسشنامه \lr{HTN-SCP} به‌صورت سه‌گانه شناخت، رفتار و عملکرد قابل سنجش است \cite{riegel2012middlerange,branch2009addie,han2014hscp}.
\par
این چارچوب مفهومی، پژوهش حاضر را بر یک بنیاد نظری استوار قرار می‌دهد و امکان تحلیل روش‌مند و دقیق اثربخشی نسبی روش‌های آموزشی را فراهم می‌سازد. مداخله آموزشی به‌عنوان متغیر مستقل تلقی می‌شود که با اثر بر متغیرهای میانجی (دانش، خودکارآمدی و انگیزه)، بر پیشاوردهای وابسته (عملکرد خودمراقبتی و فشار خون) اثر می‌گذارد. مقایسه کارگاهی سنتی با کارگاهی همراه با خرده‌یادگیری امکان استخراج ارزش افزوده خرده‌یادگیری را فراهم می‌آورد و به رفع خلأ دانشی موجود در این حوزه پاسخ می‌دهد.

\begin{figure}[H]
  \centering
  \begin{minipage}{0.88\textwidth}
    \centering
    \fbox{\parbox{0.76\textwidth}{\centering
      آموزش کارگاهی مشترک در هر دو گروه\\
      \textbf{افزودن خرده‌یادگیری ویدئویی در گروه مداخله}}}\\[0.45em]
    $\Downarrow$\\[0.45em]
    \fbox{\parbox{0.76\textwidth}{\centering
      تکرار و دسترسی مستمر به پیام‌های آموزشی\\
      ارتقای اطلاعات، آگاهی، خودکارآمدی و انگیزه}}\\[0.45em]
    $\Downarrow$\\[0.45em]
    \fbox{\parbox{0.76\textwidth}{\centering
      بهبود عملکرد خودمراقبتی\\
      تغییر فشار خون سیستولیک و دیاستولیک}}
  \end{minipage}
  \caption{چارچوب مفهومی پژوهش: ارتباط روش آموزشی با عملکرد خودمراقبتی}
  \label{fig:conceptual-framework}
\end{figure}

\section{مروری بر مطالعات پیشین}
\label{sec:previous-studies}
مرور مطالعات انجام‌شده درباره آموزش خودمراقبتی در افراد مبتلا به پرفشاری خون نشان می‌دهد که مداخلات آموزشی، با وجود تفاوت در الگوی نظری، شیوه ارائه، مدت پیگیری و ابزار سنجش، اغلب با بهبود دانش، نگرش، پایبندی به درمان و برخی رفتارهای خودمراقبتی همراه بوده‌اند. بااین‌حال، میزان اثر این مداخلات بر پیامدهای عینی مانند فشار خون و نیز پایداری تغییرات رفتاری در طول زمان یکسان نبوده است. بخش عمده پژوهش‌ها اثربخشی یک روش آموزشی را در مقایسه با مراقبت معمول سنجیده‌اند و مطالعاتی که ارزش افزوده یک جزء آموزشی کوتاه، تکرارشونده و دیجیتال را در کنار آموزش حضوری بررسی کرده باشند، محدودند. ازاین‌رو، بررسی پیشینه پژوهش مستلزم توجه هم‌زمان به شواهد اختصاصی پرفشاری خون و شواهد مرتبط با آموزش ترکیبی و خرده‌یادگیری در سایر بیماری‌های مزمن است.

\subsection{مطالعات داخلی}
در یکی از مطالعات داخلی، یک مداخله آموزشی مبتنی بر الگوی پرسید ـ پروسید در ۱۲۰ بیمار مبتلا به پرفشاری خون بررسی شد. این الگو با تاکید بر عوامل مستعدکننده، تقویت‌کننده و تسهیل‌کننده رفتار، آموزش را از انتقال صرف اطلاعات فراتر می‌برد و موانع رفتاری و محیطی را نیز در نظر می‌گیرد. پس از اجرای مداخله، دانش و نگرش بیماران بهبود یافت و رفتارهای خودمراقبتی و کیفیت زندگی آنان ارتقا پیدا کرد. همچنین کاهش فشار خون سیستولی در گروه مداخله گزارش شد. اهمیت این مطالعه در آن است که نشان داد مداخله آموزشی ساختاریافته می‌تواند هم بر متغیرهای شناختی و هم بر یک پیامد بالینی اثر بگذارد. بااین‌حال، طراحی آن امکان پاسخ‌گویی مستقیم به این پرسش را فراهم نمی‌کرد که افزودن آموزش‌های کوتاه و تکرارشونده به جلسات حضوری تا چه اندازه بر نتیجه مداخله می‌افزاید \cite{rakhshani2024precede}.
\par
در مطالعه‌ای دیگر، آموزش خودمراقبتی متناسب با سبک یادگیری در ۴۴ بیمار مبتلا به پرفشاری خون اجرا شد. در این پژوهش، شیوه ارائه محتوا با ترجیحات یادگیری افراد هماهنگ شد و پیامد اصلی، تغییر در سبک زندگی بود. پس از مداخله، نمرات سبک زندگی در گروه‌های مورد بررسی افزایش یافت. این یافته بر اهمیت تفاوت‌های فردی در دریافت و پردازش پیام‌های آموزشی تاکید دارد و نشان می‌دهد که اثربخشی آموزش تنها به محتوای علمی وابسته نیست، بلکه تناسب قالب ارائه با ویژگی‌های یادگیرنده نیز اهمیت دارد. در عین حال، حجم نمونه محدود و تمرکز بر سبک زندگی، تعمیم نتیجه به عملکرد کلی خودمراقبتی و تغییرات فشار خون را با احتیاط روبه‌رو می‌کند \cite{ghorbani2021vark}.
\par
اثر آموزش به روش بحث گروهی نیز در یک مطالعه داخلی با مشارکت ۱۰۰ بیمار مراجعه‌کننده به مراکز بهداشتی بررسی شد. در این مداخله، تبادل تجربه، طرح پرسش، حل مسئله و مشارکت فعال بیماران جایگزین ارتباط یک‌سویه آموزشی شد. پس از آموزش، میانگین نمره رفتارهای خودمراقبتی در گروه مداخله بیشتر بود و فشار خون سیستولی و دیاستولی کاهش یافت. این نتیجه با این دیدگاه سازگار است که تعامل گروهی می‌تواند فهم عملی بیماران از توصیه‌های درمانی را افزایش دهد و از راه حمایت اجتماعی، انگیزه لازم برای تداوم رفتار را تقویت کند. بااین‌حال، بحث گروهی به حضور هم‌زمان بیماران و آموزش‌دهنده وابسته است و در فاصله میان جلسات، امکان مرور منظم و استاندارد محتوا را به‌تنهایی فراهم نمی‌کند \cite{nasirian2019group}.
\par
در یک کارآزمایی شاهددار تصادفی‌شده با ۱۰۲ بیمار، مداخله‌ای مبتنی بر الگوی پرسید برای ارتقای رفتارهای خودمراقبتی اجرا شد. پس از مداخله، نمره سازه‌های الگو و رفتارهای خودمراقبتی افزایش یافت. تصادفی‌سازی و وجود گروه شاهد از نقاط قوت این پژوهش بود و احتمال انتساب تغییرات مشاهده‌شده به مداخله را افزایش می‌داد. با وجود این، تمرکز اصلی مطالعه بر اثربخشی یک برنامه مبتنی بر الگوی رفتاری بود و مقایسه میان دو قالب فعال آموزش، یعنی آموزش حضوری و آموزش حضوری تقویت‌شده با محتوای کوتاه دیجیتال، در آن انجام نشد \cite{abedi2020precede}.
\par
در پژوهشی دیگر، آموزش گروهی برای ۱۱۲ زن در سنین ۲۰ تا ۴۹ سال با هدف پیشگیری و کنترل پرفشاری خون بررسی شد. پس از مداخله، آگاهی، نگرش و رفتارهای پیشگیرانه در گروه آموزش‌دیده بهبود یافت. این مطالعه نشان داد که بحث گروهی می‌تواند حتی در جمعیتی که الزاما همگی بیمار مبتلا به پرفشاری خون نیستند، به اصلاح عوامل شناختی و رفتاری مرتبط با بیماری کمک کند. بااین‌حال، تفاوت جامعه هدف با بیماران دارای تشخیص قطعی، کاربرد مستقیم یافته‌ها را برای خودمراقبتی درمانی محدود می‌سازد؛ زیرا رفتارهای پیشگیرانه با اقداماتی مانند مصرف منظم دارو، پایش خانگی فشار خون و پیگیری درمان یکسان نیستند \cite{ramazankhani2017group}.
\par
مجموع مطالعات داخلی نشان می‌دهد که آموزش ساختاریافته، آموزش مبتنی بر الگوهای رفتاری، تناسب آموزش با سبک یادگیری و مشارکت فعال در بحث گروهی می‌توانند رفتارهای مرتبط با خودمراقبتی را بهبود دهند. جهت بیشتر یافته‌ها مثبت است و در برخی مطالعات، کاهش فشار خون نیز گزارش شده است. بااین‌حال، ناهمگونی در حجم نمونه، جامعه مورد بررسی، محتوای آموزشی، طول دوره مداخله، مدت پیگیری و نوع پیامد، مقایسه مستقیم نتایج را دشوار می‌کند. افزون بر این، بیشتر پژوهش‌ها یک مداخله آموزشی را با مراقبت معمول یا وضعیت پیش از مداخله مقایسه کرده‌اند. در نتیجه، هنوز روشن نیست که آیا افزودن خرده‌یادگیری ویدئویی به یک برنامه کارگاهی، نسبت به همان برنامه کارگاهی به‌تنهایی، به بهبود بیشتری در عملکرد خودمراقبتی بیماران مبتلا به پرفشاری خون منجر می‌شود یا خیر.

\subsection{مطالعات خارجی}
در یک مطالعه با ۶۰ بیمار مبتلا به پرفشاری خون در یک منطقه روستایی، آموزش از راه دور پرستاری بر پایه مراقبت حمایتی اجرا شد. ارتباط از راه دور امکان می‌داد آموزش و پیگیری بیماران بدون وابستگی کامل به حضور فیزیکی ادامه یابد. پس از مداخله، کیفیت زندگی بیماران گروه مداخله بهبود پیدا کرد. این نتیجه نشان می‌دهد که تداوم تماس آموزشی و حمایتی می‌تواند محدودیت دسترسی جغرافیایی را کاهش دهد و بیماران را در مدیریت روزمره بیماری یاری کند. بااین‌حال، کیفیت زندگی پیامدی چندبعدی است و بهبود آن لزوما به معنای تغییر هم‌زمان همه رفتارهای خودمراقبتی یا کنترل بهتر فشار خون نیست \cite{suprayitno2023telenursing}.
\par
در یک کارآزمایی تصادفی‌شده با ۱۳۷ بیمار، برنامه حمایت از خودمدیریتی از نظر دانش، پایبندی به درمان، مدیریت خودمراقبتی و فشار خون بررسی شد. یافته‌ها افزایش دانش و پایبندی به درمان و نیز کاهش فشار خون سیستولی و دیاستولی را نشان دادند. این مطالعه از آن جهت مهم است که در کنار پیامدهای ذهنی و پرسشنامه‌ای، تغییر شاخص بالینی را نیز ارزیابی کرد. نتایج آن بیانگر این است که آموزش زمانی می‌تواند اثر بالینی بیشتری داشته باشد که با حمایت از تصمیم‌گیری، پیگیری رفتار و تقویت نقش فعال بیمار همراه شود. بااین‌حال، برنامه حمایت از خودمدیریتی مجموعه‌ای از اجزای مداخله‌ای بود و سهم مستقل قالب آموزشی کوتاه و تکرارشونده از سایر اجزا مشخص نبود \cite{kurt2022selfmgmt}.
\par
فناوری آموزشی مبتنی بر فلش‌کارت در مطالعه‌ای با ۱۱۶ بیمار مبتلا به پرفشاری خون ارزیابی شد. محتوای فشرده و قطعه‌بندی‌شده فلش‌کارت‌ها امکان می‌داد پیام‌های کلیدی به شکلی ساده و قابل مرور در اختیار بیمار قرار گیرد. پس از مداخله، کیفیت زندگی و پایبندی به درمان بهبود یافت. این یافته از قابلیت ابزارهای آموزشی مختصر برای تقویت یادآوری و دسترسی مکرر به توصیه‌ها حمایت می‌کند. بااین‌حال، فلش‌کارت با خرده‌یادگیری ویدئویی یکسان نیست؛ زیرا ویدئو می‌تواند نمایش عملی مهارت، ترکیب تصویر و صدا و هدایت گام‌به‌گام بیمار را فراهم کند. بنابراین، نتایج این مطالعه تنها به‌صورت غیرمستقیم از استفاده از محتوای کوتاه دیجیتال پشتیبانی می‌کند \cite{desouza2016flashcard}.
\par
برآیند مطالعات خارجی اختصاصی پرفشاری خون نشان می‌دهد که آموزش از راه دور، حمایت از خودمدیریتی و فناوری‌های آموزشی ساده می‌توانند بر کیفیت زندگی، پایبندی درمانی، دانش و گاه فشار خون اثر مطلوب داشته باشند. وجه مشترک این مداخلات، عبور از آموزش مقطعی و ایجاد فرصت برای تکرار، حمایت یا دسترسی آسان‌تر به محتواست. با وجود این، تفاوت در اجزای مداخله و پیامدهای سنجیده‌شده مانع از آن است که برتری یک قالب خاص با قطعیت نتیجه‌گیری شود.


\subsection{شواهد آموزش کارگاهی در حوزه خودمراقبتی}
آموزش کارگاهی بر مشارکت فعال یادگیرنده، تمرین عملی، گفت‌وگوی گروهی و بازخورد فوری استوار است. در مطالعات انجام‌شده در گروه‌های مختلف، از بیماران مبتلا به بیماری‌های مزمن تا کارکنان سلامت و فراگیران حرفه‌ای، کارگاه‌های تعاملی با بهبود مهارت‌های خودمراقبتی، خودکارآمدی، رضایت از خودمراقبتی و آمادگی برای تغییر رفتار همراه بوده‌اند. ویژگی مهم کارگاه، تبدیل مفاهیم انتزاعی به تکلیف و مهارت عملی است. برای بیمار مبتلا به پرفشاری خون، این ویژگی می‌تواند در آموزش اندازه‌گیری صحیح فشار خون، تنظیم مصرف نمک، فعالیت بدنی، مصرف دارو و شناسایی نشانه‌های هشدار سودمند باشد.
\par
بااین‌حال، کارگاه حضوری با محدودیت‌هایی نیز همراه است. یادگیری در یک بازه زمانی محدود رخ می‌دهد و پس از پایان جلسه، بخشی از اطلاعات ممکن است فراموش شود. تراکم محتوا، تفاوت سرعت یادگیری بیماران و نبود امکان بازبینی یکسان برای همه افراد می‌تواند اثربخشی کارگاه را کاهش دهد. همچنین حضور در زمان و مکان معین برای برخی بیماران دشوار است. بنابراین، هرچند کارگاه بستر مناسبی برای تعامل و تمرین فراهم می‌کند، برای حفظ آموخته‌ها و انتقال آن‌ها به رفتار روزمره ممکن است به تقویت‌کننده‌های آموزشی پس از جلسه نیاز داشته باشد.
\par
شواهد موجود نشان می‌دهد که جو حمایتی، فرصت پرسش، تمرین عملی و بازخورد آموزش‌دهنده از سازوکارهای اصلی اثرگذاری کارگاه هستند. در مقابل، کارگاه‌هایی که عمدتا به سخنرانی فشرده محدود می‌شوند، احتمال کمتری دارند که به تغییر پایدار رفتار منجر شوند. ازاین‌رو، در تفسیر نتایج مطالعات کارگاهی باید میان حضور صرف در جلسه و مشارکت فعال در فرایند یادگیری تفاوت قائل شد.

\subsection{شواهد خرده‌یادگیری و آموزش ترکیبی}
خرده‌یادگیری محتوای آموزشی را به واحدهای کوتاه، متمرکز و مستقل تقسیم می‌کند. هر واحد معمولا یک هدف محدود دارد و می‌تواند در زمان کوتاه مشاهده یا مرور شود. کوتاهی محتوا بار شناختی را کاهش می‌دهد، امکان تکرار را افزایش می‌دهد و استفاده از آموزش را در برنامه روزمره بیمار آسان‌تر می‌سازد. در قالب ویدئویی، بیمار می‌تواند یک رفتار مشخص را مشاهده کند؛ برای نمونه، نحوه صحیح قرار دادن بازوبند دستگاه فشار خون یا انتخاب مواد غذایی کم‌نمک را چند بار ببیند و در زمان نیاز به آن بازگردد.
\par
یک مرور نظام‌مند درباره خرده‌یادگیری نشان داد که این رویکرد در بیشتر مطالعات به بهبود پیامدهای شناختی منجر شده است، اما تغییر رفتار با ثبات و فراوانی کمتری مشاهده شده بود. این تفاوت اهمیت اساسی دارد؛ زیرا دانستن توصیه‌های خودمراقبتی شرط لازم تغییر رفتار است، اما به‌تنهایی کافی نیست. برای تبدیل دانش به عمل، بیمار باید انگیزه، خودکارآمدی، فرصت تمرین، بازخورد و حمایت لازم را نیز داشته باشد. همچنین بخش قابل توجهی از مطالعات خرده‌یادگیری با محدودیت‌های روش‌شناختی و پیگیری کوتاه‌مدت روبه‌رو بوده‌اند. در نتیجه، شواهد موجود بیشتر از نقش خرده‌یادگیری به‌عنوان مکمل آموزش حمایت می‌کند تا جایگزین کامل تعامل حضوری \cite{wang2020microlearning}.
\par
در مطالعات بیماران مبتلا به دیابت نوع دو، آموزش ترکیبی شامل کارگاه و محتوای دیجیتال با افزایش دانش و رفتارهای خودمراقبتی همراه بوده است. در یک مطالعه، کارگاه حضوری همراه با آموزش اینترنتی در مقایسه با آموزش معمول، پس از دو ماه موجب بهبود ابعاد خودمراقبتی شد. در مطالعه‌ای دیگر، خرده‌یادگیری ارائه‌شده از طریق رسانه اجتماعی در مقایسه با یک کارگاه متعارف، طی پیگیری کوتاه‌مدت به بهبود بیشتر دانش، خودکارآمدی و رفتارهای خودمراقبتی انجامید. اگرچه دیابت و پرفشاری خون از نظر رژیم درمانی و پایش بالینی یکسان نیستند، هر دو بیماری به تصمیم‌گیری روزانه، پایبندی دارویی، اصلاح سبک زندگی و پیگیری طولانی‌مدت نیاز دارند. ازاین‌رو، این یافته‌ها قابلیت انتقال مفهومی به آموزش خودمراقبتی در پرفشاری خون دارند، اما برای نتیجه‌گیری قطعی باید در جمعیت هدف آزمون شوند \cite{ghiyasvandian2023blended,rahbar2024socialmedia}.
\par
در سالمندان نیز آموزش مبتنی بر خرده‌یادگیری با بهبود سواد سلامت و پایبندی دارویی همراه بوده است. این یافته برای پرفشاری خون اهمیت دارد؛ زیرا شیوع بیماری با افزایش سن بیشتر می‌شود و فراموشی توصیه‌ها، تعدد داروها و دشواری پردازش محتوای طولانی می‌تواند مانع خودمراقبتی باشد. واحدهای کوتاه و تکرارشونده ممکن است یادآوری توصیه‌ها را تسهیل کنند. بااین‌حال، توانایی استفاده از ابزار دیجیتال، کیفیت بینایی و شنوایی، دسترسی به تلفن هوشمند و حمایت خانواده از عواملی هستند که می‌توانند میزان بهره‌مندی سالمندان را تعدیل کنند \cite{poorcheraghi2025microlearning}.
\par
مطالعات آموزش ترکیبی در سایر بیماری‌های مزمن نیز نشان داده‌اند که افزودن پشتیبانی دیجیتال به آموزش حضوری می‌تواند خودکارآمدی، درگیری بیمار با برنامه مراقبتی و تداوم تمرین را تقویت کند. در عین حال، شواهد فراتحلیلی درباره مداخلات خودمراقبتی در بیماری‌های مزمن، اثری مثبت اما در مجموع متوسط و ناهمگون را نشان می‌دهد. این ناهمگونی یادآور می‌شود که افزایش تعداد شیوه‌های ارائه، به‌خودی‌خود تضمین‌کننده اثربخشی بیشتر نیست. کیفیت محتوا، تناسب آن با نیاز بیمار، میزان مشارکت، تکرار هدفمند و پیوند آموزش با رفتارهای قابل اجرا، عواملی تعیین‌کننده‌تر از تنوع ظاهری رسانه‌ها هستند \cite{munro2018blended,jennatfereidooni2026orem,lee2022chronicmeta}.

\subsection{عوامل موثر بر پیامد مداخلات آموزشی}
پیامد مداخله آموزشی به ویژگی‌های بیمار، طراحی محتوا و شرایط اجرا وابسته است. سطح اولیه دانش و خودمراقبتی، سواد سلامت، سن، وضعیت اقتصادی و اجتماعی، دسترسی به فناوری، انگیزه، حمایت خانواده و شدت بیماری می‌توانند پاسخ فرد به آموزش را تغییر دهند. افرادی که در آغاز نیاز آموزشی بیشتری دارند، گاه بهبود بیشتری نشان می‌دهند؛ اما اگر موانع دسترسی یا توانایی استفاده از فناوری برطرف نشود، ممکن است از آموزش دیجیتال کمتر بهره‌مند شوند.
\par
در سطح مداخله، محتوای اختصاصی و قابل اجرا، تکرار منظم، تعامل با آموزش‌دهنده، تمرین مهارت و بازخورد از عوامل تقویت‌کننده اثر هستند. خرده‌یادگیری زمانی احتمال بیشتری دارد که به تغییر رفتار منجر شود که با برنامه کارگاه هماهنگ باشد و هر ویدئو یک پیام روشن و یک اقدام مشخص را تقویت کند. ارسال انبوه مطالب پراکنده یا تکرار بدون هدف ممکن است موجب خستگی و کاهش مشارکت شود. همچنین، پیامدهای نزدیک مانند دانش و خودکارآمدی معمولا زودتر از پیامدهای رفتاری و بالینی تغییر می‌کنند. ازاین‌رو، ارزیابی هم‌زمان عملکرد خودمراقبتی و فشار خون برای قضاوت جامع‌تر درباره اثربخشی مداخله ضرورت دارد.
\par
مدت پیگیری نیز در تفسیر نتایج اهمیت دارد. بهبود بلافاصله پس از آموزش می‌تواند بازتاب یادآوری کوتاه‌مدت باشد، درحالی‌که سنجش پس از گذشت زمان نشان می‌دهد آیا بیمار توانسته است آموخته‌ها را در زندگی روزمره به کار گیرد یا خیر. بسیاری از مطالعات پیشین پیگیری کوتاه داشته‌اند و درباره پایداری رفتار اطلاعات محدودی ارائه کرده‌اند. در پژوهش‌های مقایسه‌ای، یکسان‌بودن محتوای اصلی کارگاه و تفاوت‌داشتن تنها جزء تقویتی خرده‌یادگیری می‌تواند به شناسایی اثر افزوده این روش کمک کند.

\subsection{جمع‌بندی مرور متون و شکاف پژوهش}
شواهد داخلی و خارجی در مجموع از اثربخشی آموزش خودمراقبتی برای بیماران مبتلا به پرفشاری خون حمایت می‌کنند. برنامه‌های مبتنی بر الگوهای رفتاری، بحث گروهی، حمایت از خودمدیریتی، آموزش از راه دور و فناوری‌های آموزشی کوتاه توانسته‌اند دانش، نگرش، سبک زندگی، پایبندی درمانی، کیفیت زندگی و رفتارهای خودمراقبتی را به درجات مختلف بهبود دهند. در برخی مطالعات، کاهش فشار خون نیز مشاهده شده است. بااین‌حال، پیامدهای شناختی به‌طور معمول سریع‌تر و باثبات‌تر از تغییرات رفتاری گزارش شده‌اند و شواهد مربوط به دوام اثر و پیامدهای عینی همچنان محدود است.
\par
آموزش کارگاهی مزیت تعامل، تمرین و بازخورد مستقیم را دارد، اما از نظر زمان، مکان و امکان مرور محتوا محدود است. خرده‌یادگیری دسترسی، تکرار و تمرکز بر پیام‌های کوچک را تسهیل می‌کند، ولی در صورت استفاده مستقل ممکن است برای ایجاد تغییر رفتاری پایدار کافی نباشد. ترکیب این دو روش از نظر نظری می‌تواند نقاط قوت هر دو را در کنار هم قرار دهد؛ کارگاه برای شکل‌گیری فهم، انگیزه و مهارت اولیه و ویدئوهای کوتاه برای مرور، یادآوری و تقویت رفتار در فاصله میان جلسات.
\par
با وجود این منطق و نتایج امیدوارکننده در برخی بیماری‌های مزمن، شواهد مستقیم درباره افراد مبتلا به پرفشاری خون محدود است. به‌ویژه، مطالعه‌ای که در آن محتوای کارگاهی دو گروه یکسان باشد و اثر افزوده خرده‌یادگیری ویدئویی بر عملکرد خودمراقبتی و فشار خون به‌طور مستقیم سنجیده شود، در پیشینه داخلی گزارش نشده است. این کمبود، تشخیص برتری واقعی روش ترکیبی را دشوار می‌سازد. پژوهش حاضر با مقایسه آموزش خودمراقبتی به روش کارگاهی و آموزش کارگاهی تلفیق‌شده با خرده‌یادگیری، در پی پاسخ به این شکاف است. سنجش هم‌زمان نمره خودمراقبتی و فشار خون می‌تواند روشن سازد که آیا تقویت دیجیتال آموزش حضوری تنها به بهبود ادراک و دانش منجر می‌شود یا قابلیت تبدیل‌شدن به رفتار خودمراقبتی و پیامد بالینی را نیز دارد.
